\documentclass[11pt]{article}
\usepackage[utf8]{inputenc}
\usepackage[T1]{fontenc}
\usepackage{geometry}
\geometry{a4paper, margin=1in}
\usepackage{titlesec}
\usepackage{enumitem}
\usepackage{booktabs}
\usepackage{array}
\usepackage{tabularx}
\usepackage{ragged2e}
\usepackage{setspace}
\onehalfspacing

\titleformat{\section}{\large\bfseries}{\thesection}{1em}{}
\titleformat{\subsection}{\normalsize\bfseries}{\thesubsection}{1em}{}

\begin{document}

\begin{center}
    \textbf{\Large Research Proposal}
\end{center}

\vspace{0.4cm}

\noindent \textbf{Samantha Irving} \quad samantha.irving@stud.unilu.ch \hfill \textbf{Date:} 01.05.2026

\vspace{0.4cm}

% ============================================================
\noindent \textbf{Preliminary title:} \\
Coercion and Hunger: District-Level Land Use Consolidation under Rwanda's
Crop Intensification Program, Household Food Welfare, and Climate
Vulnerability.

\vspace{0.4cm}

% ============================================================
\noindent \textbf{Research question:} \\
How does district-level intensity of land use consolidation under Rwanda's
Crop Intensification Program affect household food consumption welfare, and
does this effect vary by household wealth?

\vspace{0.3cm}
\noindent \textbf{Sub-questions:}
\begin{itemize}[noitemsep, topsep=2pt]
    \item Is district-level LUC intensity negatively associated with household
          food consumption per adult equivalent?
    \item Does this effect increase as household wealth decreases?
\end{itemize}

\vspace{0.4cm}

% ============================================================
\noindent \textbf{Current knowledge:} \\

Rwanda's Crop Intensification Program (CIP), launched in 2007 and ongoing,
mandates smallholder farmers to cultivate state-assigned monocrops on
consolidated plots, enforced through the \textit{imihigo} performance
contract system (Cioffo et al., 2016). The programme is widely cited as a
success, producing nationally reported yield increases---though these have
been contested on data quality grounds (Heinen, 2022). Existing research
reveals a distributional cost: Del~Prete et al.\ (2019), using EICV panel
data from 2010/11 and 2013/14, find that land consolidation participation
significantly reduced consumption of meat, fish, and fruit. Clay and Zimmerer
(2020) provide rich quantitative and qualitative evidence from four communities
in Nyamagabe district showing that CIP fields experience more than twice the
rate of climate-related yield losses compared to non-CIP fields, and that
poorer households show the greatest declines in coping capacity. However,
their findings are geographically confined to a single district in southwest
Rwanda and have not been tested at national scale using representative survey
data. This study addresses that gap, using nationally representative household
data matched to district-level LUC intensity to test whether CIP's negative
welfare effects are observable across Rwanda as a whole. The broader relevance
is clear: food insecurity driven by structural agricultural policy is an
established pathway from climate stress to household-level instability
(Hsiang et al., 2013).

\vspace{0.3cm}
The core argument is that \textit{imihigo} enforcement removed the micro-level
adaptive strategies smallholders use to manage risk---intercropping, crop
diversification across moisture gradients, and flexible plot allocation---
creating structural vulnerability (Clay and Zimmerer, 2020; Matin et al.,
2018). CIP households locked into monocrops of government-selected crops,
which Clay and Zimmerer show are among the riskiest in drought and heavy
rain, have fewer options to absorb losses when shocks hit. This cost falls
hardest on the poorest households, who additionally lack savings, credit
access, and non-farm income as alternative buffers (Matin et al., 2018).

\vspace{0.3cm}
\noindent \textbf{Hypotheses:}
\begin{itemize}[noitemsep, topsep=2pt]
    \item[\textbf{H1}] Households in higher-LUC districts will have lower
          food consumption per adult equivalent than those in lower-LUC
          districts (\textit{food welfare effect}).
    \item[\textbf{H2}] This effect will be larger for poorer households
          (\textit{poverty-concentration hypothesis}).
\end{itemize}

\vspace{0.4cm}

% ============================================================
\noindent \textbf{Preliminary structure:}
\begin{enumerate}[noitemsep, topsep=2pt]
    \item \textit{From Policy Success to Distributional Question} ---
          Introduction and relevance
    \item \textit{What the Literature Misses} ---
          State of the art and the untested resilience claim
    \item \textit{Coercion, Specialisation, and Structural Vulnerability} ---
          Theory and hypotheses
    \item \textit{Measuring the Cost} --- Research design and operationalization
    \begin{enumerate}[label=4.\arabic*, noitemsep]
        \item Data and variables
        \item Methodology and results
    \end{enumerate}
    \item \textit{Who Bears the Burden?} --- Conclusion and implications
\end{enumerate}

\vspace{0.4cm}

% ============================================================
\noindent \textbf{Research design:} \\

The study matches district-level LUC intensity---constructed as a
population-weighted average share of agricultural land under consolidation
from the NISR Seasonal Agricultural Survey (SAS) 2024 across all three
seasons---to household-level data from the EICV7 (2023/24) survey
($n \approx 15{,}000$). The dependent variable is log food consumption per
adult equivalent, drawn from the EICV7 poverty file. The analysis uses OLS
regression with LUC intensity as the main independent variable, a LUC
$\times$ wealth-quintile interaction term to test H2, and province fixed
effects and urban/rural controls throughout. The climate and conflict
relevance is addressed in the theoretical framework: CIP forces households
into drought-vulnerable monocrops precisely in districts where agricultural
climate shocks are most prevalent, and reduced food welfare is an established
antecedent of household-level instability.

\vspace{0.4cm}

% ============================================================
\noindent \textbf{Data sources:}

\vspace{0.3cm}
\noindent
\begin{tabularx}{\textwidth}{>{\bfseries}p{3.0cm} >{\RaggedRight}X >{\RaggedRight}p{3.6cm}}
\toprule
\normalfont\textbf{Variable} & \textbf{Source} & \textbf{Access} \\
\midrule
Food consumption per adult equivalent
    & EICV7 (2023/24) poverty file
    & World Bank Microdata Catalog; free registration \\
\hline
Household controls
    & EICV7 poverty file; savings \& credit modules
    & Included in EICV7 package \\
\hline
LUC intensity
    & NISR Seasonal Agricultural Survey 2024, Seasons A, B \& C
    & NISR microdata catalog; free registration \\
\bottomrule
\end{tabularx}

\vspace{0.4cm}

% ============================================================
\noindent \textbf{Literature:} \\

\noindent Cantore, Nicola. n.d. \textit{The Crop Intensification Program in
Rwanda: A Sustainability Analysis}. London: Overseas Development Institute.
\vspace{0.2cm}

\noindent Cioffo, Giuseppe, An Ansoms, and Jude Murison. 2016.
``Modernising Agriculture through a `New' Green Revolution.''
\textit{Review of African Political Economy} 43(148): 188--204.
\vspace{0.2cm}

\noindent Clay, Nathan, and Karl S. Zimmerer. 2020. ``Who Is Resilient in
Africa's Green Revolution?'' \textit{Land Use Policy} 97: 104558.
\vspace{0.2cm}

\noindent Del~Prete, Davide, Luc Ghins, Emiliano Magrini, and Karl Pauw.
2019. ``Land Consolidation, Specialization and Household Diets: Evidence
from Rwanda.'' \textit{Food Policy} 83: 139--149.
\vspace{0.2cm}

\noindent Heinen, Sebastian. 2022. ``Rwanda's Agricultural Transformation
Revisited.'' \textit{Journal of Development Studies} 58(10): 2044--2064.
\vspace{0.2cm}

\noindent Hsiang, Solomon M., Marshall Burke, and Edward Miguel. 2013.
``Quantifying the Influence of Climate on Human Conflict.''
\textit{Science} 341(6151): 1235367.
\vspace{0.2cm}

\noindent Ingelaere, Bert. 2014. ``What's on a Peasant's Mind?''
\textit{Journal of Eastern African Studies} 8(2): 214--230.
\vspace{0.2cm}

\noindent Matin, Nilufar, John Forrester, and Jonathan Ensor. 2018.
``What Is Equitable Resilience?'' \textit{World Development} 109: 197--205.
\vspace{0.2cm}

\noindent Nilsson, P\"{a}r. 2018. ``The Role of Land Use Consolidation in
Improving Crop Yields among Farm Households in Rwanda.''
\textit{Journal of Development Studies} 55(1): 1--17.
\vspace{0.2cm}

\noindent Nsabimana, Alain, et al. 2023. ``The Short and Long Run Effects of
Land Use Consolidation Programme on Farm Input Uptakes.''
\textit{Land Use Policy} 132: 106787.

\end{document}